\documentclass[dvips]{jreport}    
%\documentclass{ujarticle} 
\usepackage{amsmath,amssymb} 
%\documentclass[11pt,dvips,lines=50]{jsarticle}
\usepackage{newtxtext,newtxmath}
%\numberwithin{equation}{section}
%\usepackage[hiresbb]{graphicx}
\usepackage{graphicx}
%\usepackage{figure}
%\usepackage[dvipdfmx]{graphicx}%エラーなしで画像が出力されなかったので{graphix}にした
\usepackage{wrapfig}
\usepackage{geometry}
\geometry{left=18mm,right=10mm,top=10mm,bottom=5mm}
\setlength{\textheight}{32cm}
%\renewcommand{\bibname}{参考文献}
\usepackage[hang,small,bf]{caption}
\usepackage[subrefformat=parens]{subcaption}
\usepackage{comment}
\usepackage{ascmac}
%\usepackage{fancybx}%入ってなかった
\usepackage{caption} %キャプションの文字サイズを変えたい
\usepackage[bold]{otf}
\usepackage{tikz-feynhand}
\usepackage[misc,geometry,weather,clock]{ifsym}
\usepackage{siunitx}
\usepackage{placeins}

\captionsetup{compatibility=false}
\def\MARU#1{{\rm\ooalign{\hfil\lower.168ex\hbox{#1}\hfil \crcr\mathhexbox20D}}}




% My commands %%%%%%%%%%%%%%%%%%%%%%%%%%%%
\newcommand{\PL}[4]{
\begin{flushleft}
{\Large\underline{\PaperPortrait \ Manual Check log (E71a Central Area)}}\\
\end{flushleft}%
\vspace{-12mm}
\begin{flushright}
{\LARGE\underline{\ ECC#1  \  PL#2 \quad 2026/\qquad / \qquad}}\\
\end{flushright}
#3\\
#4
\FloatBarrier
\newpage
}

\newcommand{\Events}[1]{
\begin{table}[h]
%\vspace{-0.7cm}
% 	\caption[R]{Scan Area}
	\begin{center}
%		\setlength{\tabcolsep}{10pt}
    			\begin{tabular}{ |p{5mm} |p{15mm}|p{26mm}|p{24mm}|p{24mm}|p{20mm}|p{40mm}|} \hline
				 &Event &film左下角に移動&オイル挿し &lens表面合わせ &移動速度調整&Comment \\
				\hline \hline
				 #1
			\end{tabular}
    		\label{R}
 	\end{center}
\end{table}%
\FloatBarrier
\newpage
}


\newcommand{\AreaList}[1]{%   コマンド定義
フィルムの位置合わせ
 \begin{table}[h]
%\vspace{-0.7cm}
% 	\caption[R]{Scan Area}
	\begin{center}
%		\setlength{\tabcolsep}{10pt}
    			\begin{tabular}{ |p{10mm}|p{22mm}|p{26mm}|p{24mm}|p{24mm}|p{20mm}|p{30mm}|} \hline
				Area&SetPlateを編集&film左下角に移動 &三か所にオイル &lens表面合わせ &移動速度調整&Comment \\
				\hline \hline
				 #1
			\end{tabular}
    		\label{R}
 	\end{center}
\end{table}%
}%

\newcommand{\Area}[1]{
	 { \centering #1}& & & & & & \rule[0mm]{0mm}{0.7mm} \\ \hline
}



\newcommand{\EventList}[1]{
スキャン(Captureをクリックすると始まる)
\begin{table}[h]
%\vspace{-0.7cm}
% 	\caption[R]{Scan Area}
	\begin{center}
%		\setlength{\tabcolsep}{10pt}
    			\begin{tabular}{ |p{5mm} |p{15mm}|p{26mm}|p{24mm}|p{24mm}|p{20mm}|p{40mm}|} \hline
				 &Event &film左下角に移動&オイル挿し &lens表面合わせ &移動速度調整&Comment \\
				\hline \hline
				 #1
			\end{tabular}
    		\label{R}
 	\end{center}
\end{table}%
}

\newcommand{\Event}[2]{
				#2 & #1 & & & & & \rule[0mm]{0mm}{0.7mm} \\ \hline
}



%%%%%%%%%%%%%%%%%%%%%%%%%%%%%%%%%%%%%%%%%%%%%


\begin{document}
\thispagestyle{empty}
\input{TrackList}


\begin{comment}
\PL{6}{025}{
	\AreaList{
		\Area{2}
		\Area{3}
		\Area{5}
		\Area{6}
	}	
}{
	\EventList{
		\Event{5573}
		\Event{634877}
		\Event{566}
		\Event{256036}
		\Event{7766}
		\Event{189993}
		\Event{10810}
		\Event{455592}
		\Event{11863}
		\Event{479279}
		\Event{13423}
		\Event{203924}
		\Event{3948}
		\Event{234257}
		\Event{7803}
		\Event{417875}
		\Event{8497}
		\Event{20992}
		\Event{9344}
		\Event{1269725}
		\Event{9456}
		\Event{188292}
		\Event{5573}
		\Event{634877}
		\Event{566}
		\Event{511182}
		\Event{7766}
		\Event{605769}
		\Event{566}
		\Event{256505}
		\Event{566}
		\Event{511182}
		
	}
}
%



\begin{table}[h]
%\vspace{-0.7cm}
% 	\caption[R]{ノイズの値}
	\begin{center}
%		\setlength{\tabcolsep}{10pt}
    			\begin{tabular}{ |p{10mm}|p{8mm}|p{12mm}|p{7mm}|p{10mm}|p{8mm}|p{8mm}|p{8mm}|p{8mm}|p{6mm}|p{6mm}|p{6mm}|p{6mm}|p{22mm}|} \hline
& & & & & \multicolumn{4}{|c|}{hole}&\multicolumn{2}{c|}{face}&\multicolumn{2}{c|}{FileCheck} &Comment \\
\cline{6-13}
Date & Time & Name & pl & Event & 1&2&3&4&1&2&{\small Graph}&Size& (errorなど) \\ \hline
\quad {\LARGE /} & &  & & & & & & &&　&&& \rule[0mm]{0mm}{10mm} \\ \hline
\quad {\LARGE /} & &  & & & & & & &　&&&& \rule[0mm]{0mm}{10mm} \\ \hline
\quad {\LARGE /} & &  & & & & & & &　&&&& \rule[0mm]{0mm}{10mm} \\ \hline
\quad {\LARGE /} & &  & & & & & & &　&&&& \rule[0mm]{0mm}{10mm} \\ \hline
\quad {\LARGE /} & &  & & & & & & &&　&&& \rule[0mm]{0mm}{10mm} \\ \hline
\quad {\LARGE /} & &  & & & & & & &&　&&& \rule[0mm]{0mm}{10mm} \\ \hline
\quad {\LARGE /} & &  & & & & & & &　&&&& \rule[0mm]{0mm}{10mm} \\ \hline
\quad {\LARGE /} & &  & & & & & & &　&&&& \rule[0mm]{0mm}{10mm} \\ \hline
\quad {\LARGE /} & &  & & & & & & &　&&&& \rule[0mm]{0mm}{10mm} \\ \hline
\quad {\LARGE /} & &  & & & & & & &　&&&& \rule[0mm]{0mm}{10mm} \\ \hline
\quad {\LARGE /} & &  & & & & & & &　&&&& \rule[0mm]{0mm}{10mm} \\ \hline
\quad {\LARGE /} & &  & & & & & & &　&&&& \rule[0mm]{0mm}{10mm} \\ \hline
\quad {\LARGE /} & &  & & & & & & &　&&&& \rule[0mm]{0mm}{10mm} \\ \hline
\quad {\LARGE /} & &  & & & & & & &　&&&& \rule[0mm]{0mm}{10mm} \\ \hline
\quad {\LARGE /} & &  & & & & & & &　&&&& \rule[0mm]{0mm}{10mm} \\ \hline
\quad {\LARGE /} & &  & & & & & & &　&&&& \rule[0mm]{0mm}{10mm} \\ \hline
\quad {\LARGE /} & &  & & & & & & &　&&&& \rule[0mm]{0mm}{10mm} \\ \hline
\quad {\LARGE /} & &  & & & & & & &　&&&& \rule[0mm]{0mm}{10mm} \\ \hline
\quad {\LARGE /} & &  & & & & & & &　&&&& \rule[0mm]{0mm}{10mm} \\ \hline
\quad {\LARGE /} & &  & & & & & & &　&&&& \rule[0mm]{0mm}{10mm} \\ \hline
\end{tabular}
    		\label{R}
 	\end{center}
\end{table}%

\end{comment}



  	  		   








\begin{comment}
%\begin{itembox}[l]{Memo}
%\vspace{7cm}
%\end{itembox}%
 \begin{figure}[h]
 \vspace{-1.0cm}
	\centering
	\begin{tikzpicture} %tikzpicture環境
		\begin{feynhand}    %feynhand環境
			\vertex (a) at (-8.5,4);
			\vertex (b) at (9,4);
			\vertex (c) at (-8.5,-6);
			\vertex (d) at (9,-6);
			\vertex (e) at (-8.2,4);
			\vertex (f) at (-5.8,4);
			\propag [plain] (a) to (e);
			\propag [plain] (f) to (b);
			\propag [plain] (b) to(d);
			\propag [plain] (c) to (d);
			\propag [plain] (c) to (a);
			\node at (-7,4) {\Large \PaperPortrait \hspace{0.1cm} {\bfseries Memo}};
		\end{feynhand}
	\end{tikzpicture}
\end{figure}
\end{comment}



%表-----------------------------------------------------------
%\begin{table}[h]
% 	\caption[R]{ノイズの値}
%	\begin{center}
%		\setlength{\tabcolsep}{10pt}
 % 		\footnotesize
%    			\begin{tabular}{ |c|c|c| } \hline
%  	  		  a & b & c   \\ \hline
%    			\end{tabular}
%    		\label{R}
% 	\end{center}
%\end{table}%
%--------------------------------------------------------------


%図-そのまま----------------------------------------------------
%\begin{figure}[h]
%    \includegraphics[keepaspectratio  hight=5cm]{1.png}
%    \caption{回路図}
%    \label{fig:1}
%\end{figure}%
%--------------------------------------------------------------
%\begin{wrapfigure} {r} {55mm}
% % \begin{center}
%    \includegraphics[bb = 0 0 247 231,width=6cm]{GM.png}
%    \caption{解の場合分け}
%    \label{fig:GM}
%  \end{center}
%\end{wrapfigure}%



%図--縦に並べる-------------------------------------------------
%\begin{figure}[t]
%  \begin{minipage}[b]{0.5\linewidth}
%    \centering
 %   \includegraphics[bb = 0 0 456 215, width=7cm]{2.png}
%    \subcaption{入力電圧と出力電圧の波形}\label{ラベル1}
%  \end{minipage}
%  \begin{minipage}[b]{0.45\linewidth}
%    \centering
%    \includegraphics[bb = 0 0 456 215, width=7cm]{3.png}
%    \subcaption{周波数に対する特性}\label{ラベル2}
%  \end{minipage}
%  \caption{波形}\label{ラベル}
%\end{figure}
%-------------------------------------------------------------
\end{document}